% new_example.tex — tests all new templates added to beamerthemepolimi
% Compile with:   xelatex new_example.tex   (recommended)
%             or  lualatex new_example.tex
%             or  pdflatex new_example.tex  (falls back to generic sans font)

\documentclass[aspectratio=169]{beamer}
\usepackage[lang=en]{beamerthemepolimi}

% ---------- Presentation metadata ----------
\title{New Template\\[0.8ex]Showcase}
\polimisubtitle{All new slides at a glance}
\author{Jane Smith}
\date{10.05.2026}
\polimiheadertext{School of Engineering / New Templates Demo}

\AtBeginSection[]{\frame{\sectionpage}}
%\AtBeginSection[]{\tocslide[currentsection]}
\AtBeginSubsection[]{\frame{\subsectionpage}}

\begin{document}

% ============================================================================
% Title slide
% ============================================================================
\begin{frame}[plain, noframenumbering]
  \titlepage
\end{frame}


% Full table of contents — useful at the start of a talk
\tocslide

% ============================================================================
% Table of contents slide  (\tocslide)
% ============================================================================
\section{Navigation}
  \setcounter{framenumber}{1}


% At each section start you can repeat the TOC highlighting only the current
% section with \tocslide[currentsection].  Here we show it manually:
\tocslide[currentsection]

\asymtextslide{TOC}
{
  On the highlighted table of contents of the previous slide...
}{
  This doesn't make sense to me if we don't have the option to deactivate the Section slide (blue one)\\[2em]
  In order to remove it just comment the line
  \texttt{\textbackslash AtBeginSection\{\textbackslash frame\{\textbackslash sectionpage\}\}}
  in the preamble of the .tex file and and uncomment the line
  \texttt{\textbackslash AtBeginSection\{\textbackslash tocslide[currentsection]\}}
}

% ============================================================================
% Content layouts
% ============================================================================
\section{Content Layouts}

% --- Three-column symmetric layout  (\threecolslide) -----------------------
\threecolslide{Three Columns}{%
  \textbf{Column A}\\[0.4em]
  First approach or concept goes here.
  \begin{itemize}
    \item Advantage one
    \item Advantage two
  \end{itemize}
}{%
  \textbf{Column B}\\[0.4em]
  Second approach or concept for comparison.
  \begin{itemize}
    \item Advantage one
    \item Advantage two
  \end{itemize}
}{%
  \textbf{Column C}\\[0.4em]
  Third approach or concept closes the comparison.
  \begin{itemize}
    \item Advantage one
    \item Advantage two
  \end{itemize}
}

% --- Full-bleed image slide  (\fullbleedslide) ------------------------------
% Pass an empty path to render the placeholder; supply a real path in practice.
% Without title overlay:
\fullbleedslide{}

% With a title overlay at the bottom (semi-transparent navy strip):
\fullbleedslide[Deep-dive: System Architecture]{}

% --- Quote / highlight slide  (\quoteslide) ---------------------------------
% Plain light background, no image
\quoteslide{``Science is the poetry of reality.''}{Richard Dawkins}

% Plain light background, no image, no attribution
\quoteslide{``Science is the poetry of reality.''}{}

% Image background, 40% desaturated
\quoteslide[image=figures/wallpaper.png, desaturate=0.6]{``Science is the poetry of reality.''}{Richard Dawkins}

% Image background with navy shade and white text
\quoteslide[image=figures/wallpaper.png, navyshade=0.55, textcolor=polimiPaper]
  {``Science is the poetry of reality.''}{Richard Dawkins}

% Navy shade only (no image), white text
\quoteslide[navyshade=0.85, textcolor=polimiPaper]{``Science is the poetry of reality.''}{Richard Dawkins}

% Heavy navy shade on image, white text, no attribution
\quoteslide[image=figures/wallpaper.png, navyshade=0.85, textcolor=polimiPaper]
  {``Science is the poetry of reality.''}{}


% --- Big number callout slide  (\bignumberslide) ----------------------------
\bignumberslide{Key Impact Metrics}{%
  78\%}{Reduction in inference latency after model distillation}{%
  123.4M}{Parameters in the compressed architecture}{%
  567 mln}{Total training tokens across all fine-tuning stages}

% ============================================================================
% Math / Academic
% ============================================================================
\section{Math and Academic}

% --- Theorem / definition slide  (\theoremslide) ----------------------------
% The block type can be Theorem, Lemma, Proposition, Definition, Corollary, …
\theoremslide{Convergence Analysis}{Theorem}{Banach Fixed-Point}{%
  Let $(X, d)$ be a non-empty complete metric space and
  $T : X \to X$ a contraction mapping with constant $q \in [0,1)$.
  Then $T$ has a unique fixed point $x^* \in X$, and for any
  $x_0 \in X$ the sequence $x_{n+1} = T(x_n)$ converges to $x^*$.
}{}

\theoremslide{Key Definition}{Definition}{Lipschitz Continuity}{%
  A function $f : \mathbb{R}^n \to \mathbb{R}^m$ is \emph{Lipschitz continuous}
  with constant $L \geq 0$ if
  \[
    \|f(x) - f(y)\| \leq L\,\|x - y\| \quad \forall\, x, y \in \mathbb{R}^n.
  \]
}{}

\theoremslide{Convergence Analysis}{Theorem}{Banach Fixed-Point}{%
  Let $(X, d)$ be a non-empty complete metric space and
  $T : X \to X$ a contraction mapping with constant $q \in [0,1)$.
  Then $T$ has a unique fixed point $x^* \in X$, and for any
  $x_0 \in X$ the sequence $x_{n+1} = T(x_n)$ converges to $x^*$.
}{
  You can also comment here on the implications of the theorem, or sketch the proof idea in a second block.
}

\twotheoremslide{Key Definition}{Definition}{Lipschitz Continuity}{%
  A function $f : \mathbb{R}^n \to \mathbb{R}^m$ is \emph{Lipschitz continuous}
  with constant $L \geq 0$ if
  \[
    \|f(x) - f(y)\| \leq L\,\|x - y\| \quad \forall\, x, y \in \mathbb{R}^n.
  \]
}{Theorem}{Banach Fixed-Point}{%
  Let $(X, d)$ be a non-empty complete metric space and
  $T : X \to X$ a contraction mapping with constant $q \in [0,1)$.
  Then $T$ has a unique fixed point $x^* \in X$, and for any
  $x_0 \in X$ the sequence $x_{n+1} = T(x_n)$ converges to $x^*$.
}


% --- Algorithm / pseudocode slide  (\algoslide) ------------------------------
% The third argument is monospace pseudocode; use \\ for line breaks and
% \textbf{keyword} for algorithm keywords.  For true verbatim code, use
% a regular frame with the listings or minted package instead.
\algoslide{Gradient Descent}{%
  A first-order iterative optimisation algorithm for finding a local
  minimum of a differentiable function.\\[0.8em]
  \textbf{Complexity:}\\
  $\mathcal{O}(n \cdot k)$ per epoch, where $n$ is the number of
  parameters and $k$ the number of samples.
}{%
  \textbf{Input:} $f$, $\nabla f$, $\alpha$, $x_0$, $T$\\
  \textbf{Output:} approximate minimiser $x^*$\\[0.4em]
  $t \leftarrow 0$\\
  \textbf{while} $t < T$ \textbf{do}\\
  \quad $g \leftarrow \nabla f(x_t)$\\
  \quad $x_{t+1} \leftarrow x_t - \alpha \cdot g$\\
  \quad $t \leftarrow t + 1$\\
  \textbf{end while}\\
  \textbf{return} $x_t$%
}

% --- Timeline slide  (\timelineslide) ----------------------------------------
\timelineslide{Project Roadmap}{%
  Q1 2025}{%
  Literature review and problem formalisation}{%
  Q2 2025}{%
  Dataset collection and baseline experiments}{%
  Q3 2025}{%
  Model training and ablation studies}{%
  Q4 2025}{%
  Paper submission and open-source release}

% ============================================================================
% Practical / Structural
% ============================================================================
\section{Practical Slides}

% --- Contact / author info slide  (\contactslide) ---------------------------
% Args: {name}{email}{institution}{photo path}{QR code path}
% Empty paths render the placeholder box.
\contactslide{%
  Prof.\ Jane Smith%
}{%
  jane.smith@polimi.it%
}{%
  Politecnico di Milano — DEIB%
}{%
  figures/genericPropic.png%
}{%
  figures/qrcode.pdf%
}

% ============================================================================
% Closing slide
% ============================================================================
\closingframe

% ============================================================================
% Appendix markers  (\appendixdivider and \backupmarker)
% ============================================================================

% \backupmarker signals "the main talk ended here" with a plain dark slide.
% Use noframenumbering (already set internally) so the total count is not
% inflated.  Place it immediately before the first backup frame.
\backupmarker

% \appendixdivider looks like a lighter section page — useful when the appendix
% covers multiple topics and you want a visual divider between them. One has to 
% insert by hand the number of the appendix in the second argument, e.g. [A], [B], etc.
\appendixdivider[An appendix]{A}

\appendixdivider[Extra Results]{B}

% --- Backup frame example ---------------------------------------------------
\begin{frame}[noframenumbering]{Supplementary: Additional Ablation}
  This frame is part of the backup section. It is not counted in the
  main slide total because \texttt{\textbackslash backupmarker} was
  inserted beforehand and this frame uses \texttt{noframenumbering}.

  \begin{block}{Tip}
    Apply \texttt{[noframenumbering]} to every backup frame so the
    header bar counter stays frozen at the last main-talk slide number.
  \end{block}
\end{frame}

\end{document}
